\section*{Bab \ref{ch:g11:organic-skeletons} --- Molekul Organik: Kerangka dan Nama}

\begin{solution}{exo:g11:organic-skeletons:1}
Keduanya memiliki enam \omterm{def:g7:atoms-and-molecules:atom}{atom} karbon (ujung dan sudut) dan \ce{C6H14}:
heksana dan 2-metilpentana adalah \omterm{def:g11:organic-skeletons:isomer}{isomer}.
\end{solution}

\begin{solution}{exo:g11:organic-skeletons:2}
Propana; heksana; 2-metilbutana.
\end{solution}

\begin{solution}{exo:g11:organic-skeletons:3}
\ce{CH3-CH2-CH3}; \chemfig{CH_3-CH(-[2]CH_3)-CH_3};
\chemfig{H-C(-[2]H)(-[6]H)-C(-[2]H)(-[6]H)-H}.
\end{solution}

\begin{solution}{exo:g11:organic-skeletons:4}
Yang pertama, 3-metilpentana: 6 \omterm{def:g7:atoms-and-molecules:atom}{atom} karbon; ujung-ujung \ce{CH3}-nya
(ada tiga) mengikat 9 \omterm{def:g7:atoms-and-molecules:atom}{atom} hidrogen, kedua sudut \ce{CH2}-nya 4, dan
titik cabang \ce{CH}-nya 1: \ce{C6H14}. Sikloheksana: 6 sudut,
masing-masing \ce{CH2}: \ce{C6H12}.
\end{solution}

\begin{solution}{exo:g11:organic-skeletons:5}
\ce{C8H18}. $2n + 2 = 20$ memberikan $n = 9$: nonana, \ce{C9H20}.
\end{solution}

\begin{solution}{exo:g11:organic-skeletons:6}
\setchemfig{atom sep=1.6em}
3-metilheksana \chemfig{-[:30]-[:-30]-[:30](-[:90])-[:-30]-[:30]-[:-30]},
\ce{C7H16}; 2,2-dimetilbutana \chemfig{-[:0](-[:90])(-[:-90])-[:0]-[:30]},
\ce{C6H14}; 3-etilpentana
\chemfig{-[:30]-[:-30](-[:-90]-[:-30])-[:30]-[:-30]}, \ce{C7H16}.
\end{solution}

\begin{solution}{exo:g11:organic-skeletons:7}
Pentana, 2-metilbutana, dan 2,2-dimetilpropana, yang digambar pada
gambar dalam bab ini.
\end{solution}

\begin{solution}{exo:g11:organic-skeletons:8}
``2-etilbutana'' memiliki rantai terpanjang lima karbon yang melewati
gugus etil: 3-metilpentana. ``4-metilpentana'' dinomori dari ujung yang
salah: 2-metilpentana. ``1-metilbutana'' sebenarnya hanyalah rantai lima
karbon: pentana.
\end{solution}

\begin{solution}{exo:g11:organic-skeletons:9}
2,2-dimetilpropana (\qty{9.5}{\celsius}) < 2-metilbutana
(\qty{27.8}{\celsius}) < pentana (\qty{36.1}{\celsius}). \omterm{def:g7:atoms-and-molecules:atom}{Atom-atom} yang
sama tersusun makin ringkas: makin sedikit sentuhan antarmolekul, makin
lemah \omterm{def:g11:polarity-and-cohesion:van-der-waals}{interaksi van der Waals}, makin rendah suhu didihnya.
\end{solution}

\begin{solution}{exo:g11:organic-skeletons:10}
Menutup rantai menjadi cincin membentuk satu ikatan \ce{C-C} lagi, yang
menggantikan dua ikatan \ce{C-H} (satu di setiap ujung). \omterm{def:g11:organic-skeletons:formulas}{Rumus
molekulnya} berbeda: bukan \omterm{def:g11:organic-skeletons:isomer}{isomer}. Sikloheksana bukan \omterm{def:g11:organic-skeletons:alkane}{alkana}
(kerangkanya berupa cincin; rumusnya bukan \ce{C_{n}H_{2n+2}}).
\end{solution}

\begin{solution}{exo:g11:organic-skeletons:11}
Butana dan 2-metilpropana (\ce{C4H10}). Propana adalah \ce{C3H8} dan
siklobutana \ce{C4H8}.
\end{solution}

\begin{solution}{exo:g11:organic-skeletons:12}
\setchemfig{atom sep=1.5em}
\begin{center}
\small
\begin{tabular}{ccc}
\chemfig{-[:30]-[:-30]-[:30]-[:-30]-[:30]} &
\chemfig{-[:30](-[:90])-[:-30]-[:30]-[:-30]} &
\chemfig{-[:30]-[:-30](-[:-90])-[:30]-[:-30]} \\
heksana & 2-metilpentana & 3-metilpentana \\[6pt]
\chemfig{-[:0](-[:90])(-[:-90])-[:0]-[:30]} &
\chemfig{-[:30](-[:90])-[:-30](-[:-90])-[:30]} & \\
2,2-dimetilbutana & 2,3-dimetilbutana &
\end{tabular}
\end{center}
\end{solution}

\begin{solution}{exo:g11:organic-skeletons:13}
Dua garis zig-zag yang panjang dapat berbaring berdampingan sepanjang
seluruh panjangnya, dengan banyak \omterm{def:g7:atoms-and-molecules:atom}{atom} yang bersentuhan; dua bola hanya
bersentuhan pada bidang yang kecil. \omterm{def:g11:polarity-and-cohesion:van-der-waals}{Interaksi van der Waals} bekerja
antara \omterm{def:g7:atoms-and-molecules:atom}{atom-atom} yang bersentuhan: interaksi itu lebih besar di antara
\omterm{def:g7:atoms-and-molecules:molecule}{molekul-molekul} rantai lurus, yang memerlukan suhu lebih tinggi untuk
dipisahkan.
\end{solution}

\begin{solution}{exo:g11:organic-skeletons:14}
Rantai terpanjang memiliki enam karbon; jika dinomori dari kiri,
substituen-substituennya berada pada 2, 4, 4, dari kanan pada 3, 3, 5;
perbedaan pertama (2 dibandingkan 3) yang menentukan:
2,4,4-trimetilheksana.
\end{solution}

\begin{solution}{exo:g11:organic-skeletons:15}
\setchemfig{atom sep=1.6em}
\chemfig{-[:0](-[:90])(-[:-90])-[:-30]-[:30](-[:90])-[:-30]}: rantai lima
karbon dengan dua gugus metil pada karbon 2 dan satu pada karbon 4,
seluruhnya delapan \omterm{def:g7:atoms-and-molecules:atom}{atom} karbon; $2 \times 8 + 2 = 18$ \omterm{def:g7:atoms-and-molecules:atom}{atom} hidrogen,
\ce{C8H18}, rumus oktana: keduanya \omterm{def:g11:organic-skeletons:isomer}{isomer}. \omterm{def:g11:organic-skeletons:formulas}{Rumus struktur ringkasnya}:
\chemfig{CH_3-C(-[2]CH_3)(-[6]CH_3)-CH_2-CH(-[2]CH_3)-CH_3}.
\end{solution}

\begin{solution}{pb:g11:organic-skeletons:1}
\setchemfig{atom sep=1.5em}
\textbf{1.} \ce{C4H10}: dengan $n = 4$, $2n + 2 = 10$.

\textbf{2.} \chemfig{H-C(-[2]H)(-[6]H)-C(-[2]H)(-[6]H)-C(-[2]H)(-[6]H)-C(-[2]H)(-[6]H)-H}.

\textbf{3.} \ce{CH3-CH2-CH2-CH3}; \chemfig{-[:30]-[:-30]-[:30]}.

\textbf{4.} \chemfig{H-C(-[2]H)(-[6]H)-C(-[6]H)(-[2,1.5]C(-[1]H)(-[3]H)(-[2]H))-C(-[2]H)(-[6]H)-H};
\chemfig{CH_3-CH(-[2]CH_3)-CH_3}; \chemfig{-[:30](-[:90])-[:-30]}.
\omterm{def:g7:atoms-and-molecules:molecule}{Molekul} itu memiliki \omterm{def:g7:atoms-and-molecules:atom}{atom-atom} yang sama, empat karbon dan sepuluh
hidrogen, yang terikat dengan urutan yang berbeda.

\textbf{5.} Ya: \omterm{def:g11:organic-skeletons:isomer}{isomer struktur} (kerangkanya berbeda).

\textbf{6.} Ketiganya: suhu didihnya di bawah \qty{20}{\celsius}.

\textbf{7.} Pada \qty{-5}{\celsius}, di bawah suhu didihnya, butana
tetap cair pada tekanan normal dan hanya melepaskan sedikit gas.

\textbf{8.} Propana (\qty{-42}{\celsius}) dan 2-metilpropana
(\qty{-11.7}{\celsius}) mendidih di bawah suhu musim dingin: keduanya
tetap memberikan banyak gas dalam \omterm{def:g4:air-a-mixture-of-gases:air}{udara} dingin.

\textbf{9.} 2-metilpropana bercabang, lebih ringkas: \omterm{def:g11:polarity-and-cohesion:van-der-waals}{interaksi van der
Waals-nya} lebih lemah.

\textbf{10.} Dengan tiga \omterm{def:g7:atoms-and-molecules:atom}{atom} karbon, satu-satunya kerangka yang mungkin
adalah rantai \ce{C-C-C}: memindahkan karbon mana pun ke tempat lain
menghasilkan rantai yang sama.

\textbf{11.} 2-metilbutana.

\textbf{12.} Rantainya harus dinomori dari ujung yang lebih dekat ke
substituen: 2, bukan 3.

\textbf{13.} 2, 3, dan 5.

\textbf{14.} Heptana, \chemfig{-[:30]-[:-30]-[:30]-[:-30]-[:30]-[:-30]}.

\textbf{15.} 2-metilheksana \chemfig{-[:30](-[:90])-[:-30]-[:30]-[:-30]-[:30]}
dan 3-metilheksana \chemfig{-[:30]-[:-30](-[:-90])-[:30]-[:-30]-[:30]}.
``4-metilheksana'' adalah 3-metilheksana yang dinomori dari ujung yang
salah.

\textbf{16.} 2,2-dimetilpentana, 2,3-dimetilpentana, 2,4-dimetilpentana,
dan 3,3-dimetilpentana.

\textbf{17.} Ya, 3-etilpentana. Pada karbon 2, gugus etil itu akan
membentuk rantai enam karbon melalui gugus itu: \omterm{def:g7:atoms-and-molecules:molecule}{molekulnya} akan menjadi
3-metilheksana.

\textbf{18.} 2,2,3-trimetilbutana,
\chemfig{-[:0](-[:90])(-[:-90])-[:0](-[:90])-[:0]}.

\textbf{19.} $1 + 2 + 4 + 1 + 1 = 9$.

\textbf{20.} Heptana memiliki 9 \omterm{def:g11:organic-skeletons:isomer}{isomer struktur}, termasuk heptana itu
sendiri.
\end{solution}
